\section{Reading limits from graphs}

Graphs are one of the best ways to understand limits, provided we read them carefully. A graph can show whether values approach one height, whether there is a hole, whether there is a jump, or whether the function grows without bound.

The graph should be read as information about nearby behavior.

\subsection*{Value versus limit on a graph}

The value \(f(a)\) is represented by the point of the graph at \(x=a\), if such a point is included. The limit \(\lim_{x\to a}f(x)\) is represented by the height approached by the graph near \(x=a\).

These may be the same, but they do not have to be the same.

\begin{examplebox}
Suppose a graph has a hole at \((2,5)\), but near \(x=2\) the curve approaches height \(5\) from both sides. Suppose also that the function has a filled point at \((2,1)\). Then
\[
f(2)=1,
\]
but
\[
\lim_{x\to2}f(x)=5.
\]
The function is not continuous at \(2\), because the value and the limit are different.
\end{examplebox}

This example is one of the most important pictures in the study of continuity. A graph can have a limit at a point even if the value at that point is different.

\subsection*{Reading one-sided limits}

To read
\[
\lim_{x\to a^-}f(x),
\]
follow the graph toward \(x=a\) from the left. To read
\[
\lim_{x\to a^+}f(x),
\]
follow the graph toward \(x=a\) from the right.

If the two approached heights agree, the two-sided limit exists. If they do not agree, the two-sided limit does not exist.

\begin{examplebox}
Suppose that as \(x\) approaches \(3\) from the left, the graph approaches height \(2\). As \(x\) approaches \(3\) from the right, the graph approaches height \(7\). Then
\[
\lim_{x\to3^-}f(x)=2,
\qquad
\lim_{x\to3^+}f(x)=7.
\]
Since these are different,
\[
\lim_{x\to3}f(x)
\]
does not exist.
\end{examplebox}

The issue is not the value at \(x=3\). The issue is that the two sides do not approach the same height.

\subsection*{Vertical and horizontal asymptotes}

If the graph grows without bound as \(x\) approaches \(a\), then there may be a vertical asymptote:
\[
x=a.
\]
For example, if
\[
\lim_{x\to a^+}f(x)=\infty,
\]
then the graph rises without bound on the right side of \(a\).

If the graph approaches a fixed height as \(x\to\infty\) or \(x\to-\infty\), there may be a horizontal asymptote:
\[
y=L.
\]
This means that far away, the graph gets close to the line \(y=L\).

\subsection*{What graphs can and cannot prove}

A graph can strongly suggest a limit, and it is an excellent tool for understanding. But a drawn graph may be approximate. If exactness is required, algebraic or theoretical reasoning may be needed.

\begin{remarkbox}
Use graphs to understand the question and predict the answer. Use algebra or definitions when an exact justification is required.
\end{remarkbox}

\begin{summarybox}
When reading limits from a graph, ask:
\begin{itemize}
\item What height is approached from the left?
\item What height is approached from the right?
\item Is the value at the point included, missing, or different from the limit?
\item Is there a hole, a jump, or a vertical asymptote?
\item What happens far to the left or far to the right?
\item Is the graph exact, or is it only a guide for reasoning?
\end{itemize}
\end{summarybox}

Reading limits from graphs is not a separate skill from computing limits. It is another language for the same idea: describing approached behavior.